% GENERATED FILE. Edit canonical lessons, then run npm run generate:books.
% canonical-source-hash: fnv1a64:b7ae4fc9526b58fb
% canonical-lessons: KA-C353-uyilu, KA-C353-varasudara, KA-C353-baluvali, KA-C353-parampare, KA-C353-pidisu

\chapter{Grandfather's Will}
\label{ch:ka-grandfather-s-will}
\hlchaptermodality{353}

\begin{chapteropening}
\textbf{By the end of this chapter:} \emph{I can say a will, an heir, a legacy, a heritage and to pester.}

Complete the last lesson of chapter 353: I can say a will, an heir, a legacy, a heritage and to pester.
\end{chapteropening}

\section[\texorpdfstring{\kn{ಉಯಿಲು} (uyilu)}{ಉಯಿಲು (uyilu)}]{\kn{ಉಯಿಲು} (uyilu) — a will (testament)}

\label{lesson:KA-C353-uyilu}



\begin{quote}
Before the new one: say the Kannada for to set free, then the Kannada for to capture.
\end{quote}

\subsection*{You'll want to know: \kn{ಉಯಿಲು}}
\textbf{\kn{ಉಯಿಲು}} — \emph{uyilu} — \textquotedblleft{}a will (testament)\textquotedblright{}.

\begin{quote}
\kn{ಅಜ್ಜ} \kn{ಸಾಯುವ} \kn{ಮುಂಚೆ} \kn{ಉಯಿಲು} \kn{ಬರೆದಿದ್ದರು}.
\end{quote}

\begin{quote}
\emph{ajja sāyuva muñce uyilu barediddaru.}
\end{quote}

\begin{quote}
Grandfather had written a will before he died.
\end{quote}

A will is \textquotedblleft{}written\textquotedblright{}: \kn{ಉಯಿಲು} \kn{ಬರೆ}.

\begin{cousinweb}
\textbf{\kn{ಉಯಿಲು}} (\emph{uyilu}) is usually traced to English \emph{will}, the legal paper, which reached the South Indian languages under British courts.
\end{cousinweb}

\begin{grammarlens}[title={What you've built}]
One more word for telling what happened.
\end{grammarlens}

\subsection*{Your turn}
\begin{itemize}
\raggedright
  \item \emph{Say it:} \emph{uyilu}
  \item \emph{Say it:} \emph{uyilu}, once more
  \item \emph{Say it:} \emph{ajja sāyuva muñce uyilu barediddaru.} — Grandfather had written a will before he died.
  \item \emph{Recall:} say the Kannada for to set free, then the Kannada for to capture, then say \emph{uyilu} again
\end{itemize}

\begin{tcolorbox}[breakable,colback=teal!4,colframe=teal!35!black,title={Before you move on}]
What does \kn{ಉಯಿಲು} mean? (\textquotedblleft{}A will (testament)\textquotedblright{}.) Say it once more.
\end{tcolorbox}
\section[\texorpdfstring{\kn{ವಾರಸುದಾರ} (vārasudāra)}{ವಾರಸುದಾರ (vārasudāra)}]{\kn{ವಾರಸುದಾರ} (vārasudāra) — an heir}

\label{lesson:KA-C353-varasudara}



\begin{quote}
Before the new one: say the Kannada for to capture, then the Kannada for a will.
\end{quote}

\subsection*{You'll want to know: \kn{ವಾರಸುದಾರ}}
\textbf{\kn{ವಾರಸುದಾರ}} — \emph{vārasudāra} — \textquotedblleft{}an heir\textquotedblright{}.

\begin{quote}
\kn{ಅಜ್ಜನ} \kn{ಮನೆಗೆ} \kn{ಅವರ} \kn{ಮೊಮ್ಮಗನೇ} \kn{ವಾರಸುದಾರ}.
\end{quote}

\begin{quote}
\emph{ajjana manege avara mommaganē vārasudāra.}
\end{quote}

\begin{quote}
The grandson is the heir to Grandfather's house.
\end{quote}

Plural: \kn{ವಾರಸುದಾರರು}, \emph{vārasudāraru}.

-\kn{ದಾರ} means \textquotedblleft{}holder\textquotedblright{}, as in \kn{ಜಮೀನುದಾರ}, a landholder.

\begin{cousinweb}
\textbf{\kn{ವಾರಸುದಾರ}} (\emph{vārasudāra}) is a loan from Arabic \emph{wāris}, \textquotedblleft{}an heir\textquotedblright{}, with the Persian ending -\emph{dār}, \textquotedblleft{}holder\textquotedblright{}; the word came through Hindustani.
\end{cousinweb}

\begin{grammarlens}[title={What you've built}]
One more word for telling what happened.
\end{grammarlens}

\subsection*{Your turn}
\begin{itemize}
\raggedright
  \item \emph{Say it:} \emph{vārasudāra}
  \item \emph{Say it:} \emph{vārasudāra}, once more
  \item \emph{Say it:} \emph{ajjana manege avara mommaganē vārasudāra.} — The grandson is the heir to Grandfather's house.
  \item \emph{Recall:} say the Kannada for to capture, then the Kannada for a will, then say \emph{vārasudāra} again
\end{itemize}

\begin{tcolorbox}[breakable,colback=teal!4,colframe=teal!35!black,title={Before you move on}]
What does \kn{ವಾರಸುದಾರ} mean? (\textquotedblleft{}An heir\textquotedblright{}.) Say it once more.
\end{tcolorbox}
\section[\texorpdfstring{\kn{ಬಳುವಳಿ} (baḷuvaḷi)}{ಬಳುವಳಿ (baḷuvaḷi)}]{\kn{ಬಳುವಳಿ} (baḷuvaḷi) — a legacy; an inheritance}

\label{lesson:KA-C353-baluvali}



\begin{quote}
Before the new one: say the Kannada for a will, then the Kannada for an heir.
\end{quote}

\subsection*{You'll want to know: \kn{ಬಳುವಳಿ}}
\textbf{\kn{ಬಳುವಳಿ}} — \emph{baḷuvaḷi} — \textquotedblleft{}a legacy; an inheritance\textquotedblright{}.

\begin{quote}
\kn{ಈ} \kn{ಹಳೆಯ} \kn{ಮನೆ} \kn{ಅಜ್ಜಿಯ} \kn{ಬಳುವಳಿ}.
\end{quote}

\begin{quote}
\emph{ī haḷeya mane ajjiya baḷuvaḷi.}
\end{quote}

\begin{quote}
This old house is Grandmother's legacy.
\end{quote}

Something handed down: property, a skill or a habit.

\begin{cousinweb}
\textbf{\kn{ಬಳುವಳಿ}} (\emph{baḷuvaḷi}) is usually traced to the native Dravidian family of \emph{baḷe}, \textquotedblleft{}to increase, to grow\textquotedblright{}; DEDR 5304 lists Kannada \emph{baḷuvaḷike}, \textquotedblleft{}increase, growth\textquotedblright{}.
\end{cousinweb}

\begin{grammarlens}[title={What you've built}]
One more word for telling what happened.
\end{grammarlens}

\subsection*{Your turn}
\begin{itemize}
\raggedright
  \item \emph{Say it:} \emph{baḷuvaḷi}
  \item \emph{Say it:} \emph{baḷuvaḷi}, once more
  \item \emph{Say it:} \emph{ī haḷeya mane ajjiya baḷuvaḷi.} — This old house is Grandmother's legacy.
  \item \emph{Recall:} say the Kannada for a will, then the Kannada for an heir, then say \emph{baḷuvaḷi} again
\end{itemize}

\begin{tcolorbox}[breakable,colback=teal!4,colframe=teal!35!black,title={Before you move on}]
What does \kn{ಬಳುವಳಿ} mean? (\textquotedblleft{}A legacy; an inheritance\textquotedblright{}.) Say it once more.
\end{tcolorbox}
\section[\texorpdfstring{\kn{ಪರಂಪರೆ} (parampare)}{ಪರಂಪರೆ (parampare)}]{\kn{ಪರಂಪರೆ} (parampare) — a heritage; a tradition}

\label{lesson:KA-C353-parampare}



\begin{quote}
Before the new one: say the Kannada for an heir, then the Kannada for a legacy.
\end{quote}

\subsection*{You'll want to know: \kn{ಪರಂಪರೆ}}
\textbf{\kn{ಪರಂಪರೆ}} — \emph{parampare} — \textquotedblleft{}a heritage; a tradition\textquotedblright{}.

\begin{quote}
\kn{ಈ} \kn{ಹಬ್ಬ} \kn{ನಮ್ಮ} \kn{ಊರಿನ} \kn{ಹಳೆಯ} \kn{ಪರಂಪರೆ}.
\end{quote}

\begin{quote}
\emph{ī habba namma ūrina haḷeya parampare.}
\end{quote}

\begin{quote}
This festival is an old tradition of our town.
\end{quote}

What passes from generation to generation: \kn{ಸಾಂಸ್ಕೃತಿಕ} \kn{ಪರಂಪರೆ}, cultural heritage.

\begin{cousinweb}
\textbf{\kn{ಪರಂಪರೆ}} (\emph{parampare}) is borrowed from Sanskrit \emph{paramparā}, \textquotedblleft{}an uninterrupted series, succession, lineage\textquotedblright{} (Monier-Williams).
\end{cousinweb}

\begin{grammarlens}[title={What you've built}]
One more word for telling what happened.
\end{grammarlens}

\subsection*{Your turn}
\begin{itemize}
\raggedright
  \item \emph{Say it:} \emph{parampare}
  \item \emph{Say it:} \emph{parampare}, once more
  \item \emph{Say it:} \emph{ī habba namma ūrina haḷeya parampare.} — This festival is an old tradition of our town.
  \item \emph{Recall:} say the Kannada for an heir, then the Kannada for a legacy, then say \emph{parampare} again
\end{itemize}

\begin{tcolorbox}[breakable,colback=teal!4,colframe=teal!35!black,title={Before you move on}]
What does \kn{ಪರಂಪರೆ} mean? (\textquotedblleft{}A heritage; a tradition\textquotedblright{}.) Say it once more.
\end{tcolorbox}
\section[\texorpdfstring{\kn{ಪೀಡಿಸು} (pīḍisu)}{ಪೀಡಿಸು (pīḍisu)}]{\kn{ಪೀಡಿಸು} (pīḍisu) — to pester; to torment}

\label{lesson:KA-C353-pidisu}



\begin{quote}
Before the new one: say the Kannada for a legacy, then the Kannada for a heritage.
\end{quote}

\subsection*{You'll want to know: \kn{ಪೀಡಿಸು}}
\textbf{\kn{ಪೀಡಿಸು}} — \emph{pīḍisu} — \textquotedblleft{}to pester; to torment\textquotedblright{}.

\begin{quote}
\kn{ಮನೆಗಾಗಿ} \kn{ಸಂಬಂಧಿಕರು} \kn{ಅಜ್ಜಿಯನ್ನು} \kn{ತುಂಬಾ} \kn{ಪೀಡಿಸಿದರು}.
\end{quote}

\begin{quote}
\emph{manegāgi sambandhikaru ajjiyannu tumbā pīḍisidaru.}
\end{quote}

\begin{quote}
The relatives pestered Grandmother a lot for the house.
\end{quote}

Past (I …): \kn{ಪೀಡಿಸಿದೆ}, \emph{pīḍiside}.

The person troubled takes -\kn{ಅನ್ನು}. Illness too \kn{ಪೀಡಿಸು}: \kn{ರೋಗ} \kn{ಪೀಡಿಸಿತು}.

\begin{cousinweb}
\textbf{\kn{ಪೀಡಿಸು}} (\emph{pīḍisu}) is built on Sanskrit \emph{pīḍā}, \textquotedblleft{}pain, damage\textquotedblright{} (CDIAL 8227), from the root \emph{pīḍ}, \textquotedblleft{}to press, to hurt\textquotedblright{}, with the Kannada verb ending -\emph{isu}.
\end{cousinweb}

\begin{grammarlens}[title={What you've built}]
One more word for telling what happened.
\end{grammarlens}

\subsection*{Your turn}
\begin{itemize}
\raggedright
  \item \emph{Say it:} \emph{pīḍisu}
  \item \emph{Say it:} \emph{pīḍisu}, once more
  \item \emph{Say it:} \emph{manegāgi sambandhikaru ajjiyannu tumbā pīḍisidaru.} — The relatives pestered Grandmother a lot for the house.
  \item \emph{Recall:} say the Kannada for a legacy, then the Kannada for a heritage, then say \emph{pīḍisu} again
\end{itemize}

\begin{tcolorbox}[breakable,colback=teal!4,colframe=teal!35!black,title={Before you move on}]
What does \kn{ಪೀಡಿಸು} mean? (\textquotedblleft{}To pester; to torment\textquotedblright{}.) Say it once more.
\end{tcolorbox}
